% =============================================================================
% glossary.tex -- plain-language definitions for uncommon terms in the monograph
% =============================================================================
\noindent This glossary gives plain definitions for terms that carry more weight
in this monograph than they usually do in casual prose. It is not a dictionary of
all the mathematics. It is a guide to the vocabulary that makes the argument
readable.

\section*{Core Vocabulary}
\addcontentsline{toc}{section}{Core vocabulary}
\markboth{Glossary}{Core vocabulary}

\begin{termnotes}
  \item[Actual.]
  The part of the ontology that has resolved into a discrete ternary state:
  $s=-1$, $0$, or $+1$. In ordinary language, it is what has happened, as
  opposed to what the substrate is poised to do.

  \item[Axiom.]
  A starting commitment of the framework. An axiom is not derived; it defines
  the model whose consequences are then studied.

  \item[Beable.]
  A quantity treated as something the substrate can actually realize, not merely
  something an outside observer can calculate. In the alpha discussion, the
  missing beable is the branch-selecting square root.

  \item[Binding law.]
  A rule that would force separate ingredients into one operational readout. For
  $\alpha$, it would have to bind the trace and determinant into the same
  rank-2 operator for an internal reason, rather than by choice.

  \item[Boundary.]
  A place where the construction stops. A boundary is not a failure by itself;
  it is useful information about what the axioms do and do not determine.

  \item[Calibration.]
  A declared bridge between dimensionless lattice units and physical units such
  as metres, seconds, or masses. Once a calibration is accepted, many dimensional
  quantities inherit it.

  \item[Closed negative.]
  A route that has been tried and ruled out within its stated scope. The result
  is preserved so the same failed route does not keep returning as if new.

  \item[Construction.]
  A step-by-step building of one object from earlier objects, like building
  integers from natural numbers or complex numbers from real pairs. In this
  monograph, construction means the reader should be able to audit how each
  object enters.

  \item[Dispositional.]
  The part of the ontology that represents tendency, capacity, or readiness to
  act. Here the flux field is dispositional: it stores what the substrate is
  poised to manifest.

  \item[Dynamical.]
  Determined by an additional rule, process, state, or measurement, rather than
  fixed by the bare structure. Saying $\alpha$ is dynamical means the current
  structural spine does not by itself select its value.

  \item[Epistemic tag.]
  A label that states what kind of claim is being made: theorem, conjecture,
  parametric insertion, open problem, and so on. The tag is part of the claim,
  not a side comment.

  \item[Foundational obstruction.]
  A named obstacle that blocks a hoped-for derivation at the level of principles,
  not merely at the level of algebraic tidying.

  \item[Parametric insertion.]
  A calculation where the formula comes from standard physics and FTD supplies
  numbers to put into it. It can be a useful cross-check, but it is not a
  derivation from the FTD axioms.

  \item[Readout.]
  The rule that turns a mathematical structure inside the framework into an
  operational quantity that could be compared with physics.

  \item[Selection.]
  A motivated choice among possibilities. It may be natural, but unless the
  framework forces it uniquely, it is not a theorem.

  \item[Structural.]
  Fixed by the architecture of the framework itself. A structural value is not
  chosen by measurement or calibration after the fact.

  \item[Substrate.]
  The underlying discrete lattice system: sites, states, local neighborhoods,
  and update rules before any continuum or physical interpretation is added.

  \item[Substrate-native.]
  Built from the substrate's own objects and rules. A substrate-native readout
  should not import the desired physical answer from outside.
\end{termnotes}

\section*{Mathematical Vocabulary}
\addcontentsline{toc}{section}{Mathematical vocabulary}
\markboth{Glossary}{Mathematical vocabulary}

\begin{termnotes}
  \item[BCC layer.]
  The body-centred-cubic layer inside the Moore neighbourhood: the eight body
  diagonal corners around a site. It is important because it carries a natural
  quarter-turn complex structure.

  \item[Branch selection.]
  Choosing one root of a quadratic rather than the other. The symmetric data of
  a quadratic know both roots together; branch selection needs the ingredient
  that distinguishes them.

  \item[CM period.]
  A special number arising from an elliptic curve with complex multiplication.
  In this monograph, $G^*$ is treated as a lemniscatic CM period appearing in
  several different forms.

  \item[Determinant.]
  For a rank-2 readout operator, the product-like invariant. In the alpha
  problem, the determinant slot is exactly where the unforced odd factor of
  $G^*$ would have to land.

  \item[Gaussian integers.]
  Numbers of the form $a+bi$ with $a$ and $b$ integers. They enter because a
  quarter-turn in a coordinate plane behaves like multiplication by $i$.

  \item[Galois-fixed.]
  Unchanged by the symmetry that swaps the two roots of the master quadratic.
  Data that are Galois-fixed can describe the pair of roots but cannot choose
  which one is the physical branch.

  \item[Lemniscatic ratio.]
  The number $G^*=\Gamma(1/4)/\Gamma(3/4)$. It is central to the construction
  and should not be confused with the related classical lemniscate constant
  $\varpi$.

  \item[Master quadratic.]
  The polynomial
  \[
    x^2-16G^{*2}x+16G^{*3}.
  \]
  Its larger root is the number that matches $\alpha^{-1}$; the polynomial
  itself is algebraic, while the physical identification remains conjectural.

  \item[Moore neighbourhood.]
  The 26 neighbours around a cubic-lattice site: face, edge, and corner
  neighbours together. It is the local causal neighbourhood of the model.

  \item[Rank-2 operator.]
  A two-dimensional linear readout object. Here it matters because such an
  operator is controlled by two invariants: trace and determinant.

  \item[Route-invariant.]
  Reached by different lines of attack in the same form. The alpha boundary is
  route-invariant because several natural attempts all reduce to the same
  missing binding obligation.

  \item[Square-root wall.]
  The obstruction represented by $\sqrt{G^*(4G^*-1)}$, the first object that
  distinguishes the two roots of the master quadratic. The current substrate
  does not force this branch-selecting object.

  \item[Trace.]
  For a rank-2 readout operator, the sum-like invariant. In the alpha problem,
  the trace is much better supported than the determinant assembly.

  \item[Zeta-regularized determinant.]
  A way to assign a determinant-like value to an infinite spectrum using zeta
  regularization. In this monograph it supplies a clean source of one factor of
  $G^*$, but not the full readout assembly.
\end{termnotes}

\section*{Phrases Used in the Argument}
\addcontentsline{toc}{section}{Phrases used in the argument}
\markboth{Glossary}{Phrases used in the argument}

\begin{termnotes}
  \item[The menu but not the dish.]
  A shorthand for the alpha boundary. The framework supplies ingredients that
  make the alpha readout plausible, but it does not force their final assembly.

  \item[Physics quarantine.]
  The separation of the mathematical construction from the physics-facing match.
  The phrase is meant to keep a good numerical match from quietly turning into a
  claimed derivation.

  \item[The spine.]
  The chain of theorem-grade algebraic results that can be studied without any
  physical interpretation.

  \item[The wall.]
  A memorable name for a mapped obstruction. It says: this is where the current
  construction stops, and this is the exact form of what is missing.
\end{termnotes}
